\section{Conclusion: Deep Causal Estimation in Context}
\label{section:conclusions}

In this primer we introduce social scientists to the emerging machine learning literature on deep learning for causal inference. To set the stage, we first provide both an intuitive introduction to fundamental deep learning concepts like representation and multi-task learning, as well as practical guidelines for training neural networks. In the main body of the article, we show how ML researchers have adapted core treatment and outcome modeling strategies to leverage the particular strengths of neural networks for heterogeneous treatment effect estimation. We follow with a discussion on inference (e.g., model selection, confidence intervals, interpretation), and closed with a prospective look at algorithms for inference from text, social networks, images, and time varying data.

Deep learning is not the only potential tool for heterogeneous treatment effect inference, and there are robust literatures exploring the usage of other methods in both the econometrics and biostatistics communities \citep{van2011targeted,chernozhukov2016double,wager2018}. While these literatures are certainly more mature, below we discuss reasons why we think the use gap between neural networks and other machine learning methods will continue to narrow, a change that we must prepare for.
 
First, neural networks are better at modeling non-linear heterogeneity (e.g., in treatment responses) than other machine learning methods. In extensive simulations, \citet{curth2021} found that when the data-generating process for treatment heterogeneity includes exponential relationships, neural networks outperformed random forests, but tree-based methods are robust when the data-generating process is built on linear functions. Neural networks were also consistently better at predicting outlier treatment effects than forests. These differences result from how the two methods model functions. While neural networks can approximate any continuous function with enough neurons, random forests must build non-linear or non-orthogonal decision boundaries using piecewise functions and average predictions. Consistent with these differences, \citet{curth2021} also find that neural networks do better when variables are constructed as continuous covariates, and vice versa when they are dichotomized.  

From a statistical perspective, the rise of semi-parametric and double machine learning frameworks has also narrowed the gap between neural networks and other types of machine learning in terms of theoretical guarantees. For example, the TMLE-inspired Dragonnet algorithm featured here is unbiased, plausibly consistent, and converges to the target estimand at a fast rate of $\frac{1}{\sqrt{n}}$. The closely-related Riezsnet double machine learning model (not featured) boasts similar guarantees \citep{chernozhukov2022riesznet}. Beyond these algorithms, there is a growing adjacent literature of model-agnostic plug-in learners (e.g., X-learner, R-learner) that can leverage the strengths of neural networks \citep{nie2021quasi,kunzel2019metalearners}.

Third, folk beliefs about the data-hungriness and uninterpretability of neural networks are overstated. Neural networks are data-hungry when over-parameterized or learning from high-dimensional data like images, but we show in the tutorials that modest-sized, well-regularized neural networks can successfully infer heterogeneous treatment effects in a naturalistic simulation of quantitative data with less than 800 units. In Section \ref{section:interpretation}, we also highlight the considerable progress in machine learning interpretability over the past five years, much of which has been on model-agnostic approaches that benefit all black-box algorithms equally.\footnote{Critics often point to out-of-bag feature importances as a particular strength of random forests, but this approach has been shown to be less accurate than model-agnostic permutation importances anyways \citep{altmann2010}.}

In our opinion, the most pressing limitation of current deep learning approaches is the difficulty of optimizing neural networks. Theoretically, this stems from a) the complexity of the loss functions which are often non-convex, and b) the ease of over-parameterizing these models to fit these functions. If neural networks are to be used as statistical estimators, statistical guarantees must be backed by optimization guarantees and/or more rigorous methods for model selection. Outside of statistical estimation, this limitation has largely been addressed through empirical testing on test data and strategic model selection. Within the statistical estimation context, this gap will likely need to be addressed by simulation-based sensitivity analyses and, in the short term, comparisons to other model families.

Moreover, there has been a lack of mature tools and empirical applications of these models. A major goal of this primer, and the tutorials in particular, is to synthesize the theoretical literature, practical training and interpretation guidelines,  and annotated code so that social scientists in one place can start using these models. Deep learning frameworks like Tensorflow and Pytorch are becoming more accessible every year, but we note that canned Python packages like Uber's causalML exist for interested readers who just want to experiment with a few of these models \citep{chen2020causalml}.
 
Despite current limitations, we believe the future of causal estimation runs through deep learning. As causal inference ventures into new settings, the flexibility of neural networks will become essential for learning from text, graph, image, video, and speech data. For time-varying settings, we believe the ability of neural networks to model non-linearities and long-range temporal dependencies will ultimately lead to solutions with net weaker assumptions than current approaches. Overall, we are optimistic and excited to see where deep causal estimation heads over the next few years.

\section{Author's Note}
The accompanying tutorials are available at \href{https://github.com/kochbj/Deep-Learning-for-Causal-Inference}{https://github.com/kochbj/Deep-Learning-for-Causal-Inference}. The tutorials use the IHDP naturalistic simulation introduced in \citet{Hill2011} as an example. The  25 covariates/features for the 747 units (139 treated) in the dataset were taken from an experiment, but Hill simulated the outcomes to create known counterfactuals. The data are available from Fredrik Johansson's website \href{https://www.fredjo.com}{https://www.fredjo.com}.